\BourbakiExercisesSection

\begin{enumerate}
\def\labelenumi{\arabic{enumi}.}
\tightlist
\item
  \begin{enumerate}
  \def\labelenumii{(\alph{enumii})}
  \tightlist
  \item
    若 A 是一个无零因子的环，证明投射 A-模的每个元素 \(\neq0\) 都是自由的。
  \end{enumerate}
\end{enumerate}

\begin{enumerate}
\def\labelenumi{(\alph{enumi})}
\setcounter{enumi}{1}
\item
  给出一个投射 A-模的例子，其零化子非零（参见 I，§ 8，第 11 号）。
\item
  证明 Z-模 Q 不是投射 Z-模。*（参见《交换代数》II，§ 5，习题 11，其中给出了整环上一个有限生成但不自由的投射模的例子。）*
\end{enumerate}

\begin{enumerate}
\def\labelenumi{\arabic{enumi}.}
\setcounter{enumi}{1}
\item
  证明每个投射 A-模都是某一族投射子模的直和，其中每个投射子模都具有可数个生成元（“卡普兰斯基定理”；将 §1 的习题 15 应用于 E 是自由模的情形）。
\item
  设 \(\mathfrak{c}\) 是一个基数且 \(\mathbf{P}\) 是一个投射 A-模，使得 \(\gamma(\mathbf{P}) \leqslant \mathfrak{c}\) (§ 1,
\end{enumerate}

练习 6）。证明若 L 是具有基数为 \(\geqslant c\) 的无限基的自由 A-模，则 P ⊕ L 与 L 同构。（注意存在一个自由 A-模 M，其基的基数为 \(\leqslant c\) ，同构于 P 与一个 A-模 Q 的直和。注意 M \(^{(N)}\) 与 P ⊕ M \(^{(N)}\) 是同构的。）

¶ 4. (a) 设 E, E' 是两个 A-模，N 是 E 的一个子模，N' 是 E' 的一个子模，使得 E 是投射的且 E/N 与 E'/N' 同构。证明存在一个正合序列

\[
0 \longrightarrow \mathrm{N} \stackrel {u} {\longrightarrow} \mathrm{E} \oplus \mathrm{N} ^ {\prime} \stackrel {v} {\longrightarrow} \mathrm{E} ^ {\prime} \longrightarrow 0.
\]

（注意存在一个同态 \(f \colon \mathbf{E} \to \mathbf{E}'\) ，使得 \(f(\mathbf{N}) \subset \mathbf{N}'\) ，它在取商后给出一个同构 \(\mathbf{E} / \mathbf{N} \to \mathbf{E}' / \mathbf{N}'\) ；以与 §1 第 7 号中正合序列（29）所用方式类似的方式定义 \(u\) 和 \(v\) 。）

\begin{enumerate}
\def\labelenumi{(\alph{enumi})}
\setcounter{enumi}{1}
\item
  从（a）推出，如果 \(\mathbf{E}'\) 是投射的， \(\mathbf{E} \oplus \mathbf{N}'\) 和 \(\mathbf{E}' \oplus \mathbf{N}\) 是同构的。
\item
  设 \(0 \to \mathbf{E}_m \to \mathbf{E}_{m-1} \to \ldots \to \mathbf{E}_1 \to \mathbf{E}_0 \to 0\) 是 A-模的一个正合序列，使得 \(\mathbf{E}_0, \mathbf{E}_1, \ldots, \mathbf{E}_{m-2}\) 是投射的。证明 A-模 \(\bigoplus_{h \geqslant 0} \mathbf{E}_{m-2h}\) 和 \(\bigoplus_{h \geqslant 0} \mathbf{E}_{m-2h+1}\) 是同构的（约定： \(\mathbf{E}_i = \{0\}\) 对于 \(i < 0\) ）.
\end{enumerate}

\begin{enumerate}
\def\labelenumi{\arabic{enumi}.}
\setcounter{enumi}{4}
\item
  设 \(u\) 是从 A-模 \(\mathbf{E}\) 到 A-模 \(\mathbf{F}\) 的线性映射， \(v = {}^t u\) 为其转置。对于每个子模 \(\mathbf{N}'\) （ \(\mathbf{F}^*\) 的）， \(\mathbf{E}\) 中正交于 \(v(\mathbf{N}')\) 的子模是 \(\overline{u}^{-1}(\mathbf{N})\) ，其中 \(\mathbf{N}\) 是 \(\mathbf{F}\) 中正交于 \(\mathbf{N}'\) 的子模.
\item
  \begin{enumerate}
  \def\labelenumii{(\alph{enumii})}
  \tightlist
  \item
    若 E 是 A-模，则 E 的每个单射自同态 u 在环 \(\operatorname{End}_{\mathbf{A}}(\mathbf{E})\) 中不是左零因子。反之，若对于 E 的每个子 A-模 \(F \neq \{0\}\) 存在 E 的一个自同态 \(v \neq 0\) ，使得 \(v(\mathbf{E}) \subset \mathbf{F}\) ，则在 \(\operatorname{End}_{\mathbf{A}}(\mathbf{E})\) 中不是左零因子的元素是单射自同态。上述条件在存在线性形式 \(x' \in E^*\) 和一个 \(x \in E\) 使得 \(\langle x, x' \rangle\) 在 A 中可逆时得到满足，特别地，当 E 是自由模时得到满足。
  \end{enumerate}
\end{enumerate}

\begin{enumerate}
\def\labelenumi{(\alph{enumi})}
\setcounter{enumi}{1}
\item
  每个自同态 \(\neq 0\) （ \(\mathbf{Z}\) -模 \(\mathbf{Q}\) 的）都是双射的，尽管不存在线性形式 \(\neq 0\) 于 \(\mathbf{Q}\) 上。
\item
  设 \(U_{p}\) 为 §1 习题 12(d) 中定义的 Z-模。证明自同态 \(x \mapsto px\) （ \(U_{p}\) 的）不是单射，并且在 \(\operatorname{End}_{\mathbf{Z}}(\mathbf{U}_{p})\) 中不是左零因子。
\end{enumerate}

\begin{enumerate}
\def\labelenumi{\arabic{enumi}.}
\setcounter{enumi}{6}
\tightlist
\item
  \begin{enumerate}
  \def\labelenumii{(\alph{enumii})}
  \tightlist
  \item
    若 \(u\) 是 A-模 \(\mathbf{E}\) 的一个满射自同态，则 \(u\) 在 \(\operatorname{End}_{\mathbf{A}}(\mathbf{E})\) 中不是右零因子。反之，若对于每个子模 \(\mathbf{F} \neq \mathbf{E}\) （ \(\mathbf{E}\) 的）存在一个线性形式 \(x^{*} \in \mathbf{E}^{*}\) ，它在 \(\mathbf{F}\) 上为零且是满射的，则 \(\operatorname{End}_{\mathbf{A}}(\mathbf{E})\) 中每个不是右零因子的元素都是满射自同态。
  \end{enumerate}
\end{enumerate}

\begin{enumerate}
\def\labelenumi{(\alph{enumi})}
\setcounter{enumi}{1}
\item
  若 E 是习题 6(c) 中定义的 Z-模，证明每个自同态 \(\neq 0\) （E 的）都是满射的，尽管在 E 上不存在线性形式 \(\neq 0\) 。
\item
  证明若 E 是自由 Z-模 \(\neq0\) ，则存在 E 的非满射自同态，它们在 \(\operatorname{End}_{\mathbf{Z}}(\mathbf{E})\) 中不是右零因子。
\end{enumerate}

\begin{enumerate}
\def\labelenumi{\arabic{enumi}.}
\setcounter{enumi}{7}
\tightlist
\item
  \begin{enumerate}
  \def\labelenumii{(\alph{enumii})}
  \tightlist
  \item
    设 E 是自由 A-模，F、G 是两个 A-模，且 \(u: \mathbf{E} \to \mathbf{G}\) , \(v: \mathbf{F} \to \mathbf{G}\) 为两个线性映射。证明若 \(u(\mathbf{E}) \subset v(\mathbf{F})\) ，则存在一个线性映射 \(w: \mathbf{E} \to \mathbf{F}\) ，使得 \(u = v \circ w\) 。
  \end{enumerate}
\end{enumerate}

\begin{enumerate}
\def\labelenumi{(\alph{enumi})}
\setcounter{enumi}{1}
\item
  给出两个非零自同态 \(u, v\) （ \(\mathbf{Z}\) -模 \(\mathbf{E}\) 的，该模定义于习题 6(c)）的例子，使得不存在自同态 \(w\) （ \(\mathbf{E}\) 的）满足 \(u = v \circ w\) （尽管 \(u(\mathbf{E}) = v(\mathbf{E}) = \mathbf{E}\) ）。
\item
  给出两个自同态 \(u, v\) （ \(\mathbf{Z}\) -模 \(\mathbf{Z}\) 的）的例子，使得 \(u^{-1}(0) = v^{-1}(0) = \{0\}\) ，但不存在自同态 \(w\) （ \(\mathbf{Z}\) 的）满足 \(u = w \circ v\) （参见 §7，习题 14）。
\end{enumerate}

\begin{enumerate}
\def\labelenumi{\arabic{enumi}.}
\setcounter{enumi}{8}
\tightlist
\item
  给定一个 A-模 E，对于 E 的每个子模 M（相应地，E* 的每个子模 M'），令 M°（相应地，M'°）表示 M 在 E* 中的正交补（相应地，M' 在 E 中的正交补）。考虑以下四个性质：
\end{enumerate}

\begin{enumerate}
\def\labelenumi{(\Alph{enumi})}
\item
  典范同态 \(c_{\mathbf{E}}: \mathbf{E} \to \mathbf{E}^{**}\) 是双射。
\item
  对于 E 的每个子模 M，典范同态 \(E^{*}/M^{\circ} \rightarrow M^{*}\) 是双射。
\item
  对于 E 的每个子模 M， \(M^{\circ\circ} = M\) .
\item
  对于 E 的每一对有序子模 M、N，
\end{enumerate}

\[
(\mathbf {M} \cap \mathbf {N}) ^ {\circ} = \mathbf {M} ^ {\circ} + \mathbf {N} ^ {\circ}.
\]

\begin{enumerate}
\def\labelenumi{(\alph{enumi})}
\item
  证明条件 (B) 蕴含 (D)（证明对于 \(x^{*} \in (\mathbf{M} \cap \mathbf{N})^{\circ}\) ，存在 \(y^{*} \in \mathbf{E}^{*}\) ，使得 \(\langle x + y, y^{*} \rangle = \langle x, x^{*} \rangle\) 对于 \(x \in \mathbf{M}\) 和 \(y \in \mathbf{N}\) 成立）。
\item
  设 A 是一个无零因子的环，但不具有左分式域 \(*\) （例如，一个维数 \textgreater1 的向量空间的张量代数，参见 III，§5，习题 5） \(^{*}\) 。若取 \(E = A_{s}\) ，则条件 (A) 成立，但条件 (B)、(C)、(D) 均不成立（关于 (B)、(C)、(D) 成立而 (A) 不成立的例子，参见 §7，第 5 号，定理 6）。
\item
  取 A 为乘积环 \(\mathbf{K}^{\mathrm{I}}\) ，其中 \(\mathbf{K}\) 是一个交换域，且 I 是无限集。证明 A-模 \(\mathbf{E} = \mathbf{A}\) 满足条件 (A) 和 (B)，但不满足 (C)（注意， \(\mathbf{A}\) 的理想的零化子都具有 \(\mathbf{K}^{\mathrm{J}}\) 的形式，其中 \(\mathbf{J} \subset \mathbf{I}\) ）。
\item
  假设对于两个子模 \(\mathbf{M},\mathbf{N}\) （ \(\mathbf{E}\) 的），满足 \(\mathbf{N} \subset \mathbf{M}\) 和 \(\mathbf{N} \neq \mathbf{M}\) ，对偶 \((\mathbf{M} / \mathbf{N})^*\) 非零；证明条件 (B) 蕴含 (C)。
\item
  \(c_{\mathbf{E}}\) 的核是正交 \((\mathbf{E}^{*})^{\circ}\) （ \(\mathbf{E}^{*}\) 在 \(\mathbf{E}\) 中的）。给出一个例子，其中 \(\mathbf{E}^{*}\) 与 \((\mathbf{E}^{*})^{\circ}\) 都不为 0（考虑一个模，它包含一个元素，该元素的零化子含有一个非零因子的元素）。
\item
  设 \(\mathbf{M}\) 是 \(\mathbf{E}\) 的一个直因子；证明 \(\mathbf{M}^{\circ \circ} = \mathbf{M} + (\mathbf{E}^{*})^{\circ}\) 。
\end{enumerate}

\begin{enumerate}
\def\labelenumi{\arabic{enumi}.}
\setcounter{enumi}{9}
\tightlist
\item
  给出一个 A-线性映射 \(u: \mathbf{E} \to \mathbf{F}\) 的例子，使得 \(^t u\) 是双射但 \(u\) 既不是单射也不是满射（参见习题 6(b) 和 (c)）。
\end{enumerate}

推导出一个例子，其中 \(y_{0} \in F\) 正交于 \(^{t}u\) 的核，但方程 \(u(x) = y_{0}\) 无解。

¶ 11. 设 A 是一个环，I 是一个 A-模。I 称为内射的，如果对每个由 A-线性映射组成的正合序列 \(\mathrm{E}' \to \mathrm{E} \to \mathrm{E}''\) ，序列

\[
\operatorname{Hom} \left(\mathrm{E} ^ {\prime \prime}, \mathrm{I}\right)\rightarrow \operatorname{Hom} (\mathrm{E}, \mathrm{I}) \rightarrow \operatorname{Hom} \left(\mathrm{E} ^ {\prime}, \mathrm{I}\right)
\]

是正合的。

\begin{enumerate}
\def\labelenumi{(\alph{enumi})}
\tightlist
\item
  证明以下性质是等价的：
\end{enumerate}

\((\alpha)\) I 是单射。

(β) 对每个由 A-线性映射组成的正合序列 \(0 \rightarrow E' \rightarrow E \rightarrow E'' \rightarrow 0\) ，序列

\[
0 \rightarrow \operatorname{Hom} (\mathrm{E} ^ {\prime \prime}, \mathrm{I}) \rightarrow \operatorname{Hom} (\mathrm{E}, \mathrm{I}) \rightarrow \operatorname{Hom} (\mathrm{E} ^ {\prime}, \mathrm{I}) \rightarrow 0
\]

是正合的。

(γ) 对每个 A-模 M 及 M 的每个子模 N，从 N 到 I 的每个 A-线性映射都能延拓为从 M 到 I 的 A-线性映射。

(δ) 对于 A 的每个左理想 \(\mathfrak{a}\) 以及每个 A-线性映射 \(f: \mathfrak{a} \to \mathbf{I}\) ，存在一个元素 \(b \in \mathbf{I}\) ，使得 \(f(a) = ab\) 对所有 \(a \in \mathfrak{a}\) 成立。

(ζ) I 是包含它的每个 A-模的一个直因子。

(θ) 对于每个 A-模 E，若 E 是 I 与一个单生成子模的和，则 I 是 E 的一个直因子。

为证明(δ)蕴含(γ)，需说明性质(δ)蕴含：若E是A-模，F是E的子模且E/F是单生成的，则F到I的每个线性映射都能延拓为E到I的线性映射。然后使用Zorn引理。为证明 (θ) 蕴含 (δ)，考虑 A-模 \(A_{s} \times I = M\) 以及 M 的由元素 \((a, -f(a))\) （对于 \(a \in \mathfrak{a}\) ）组成的子模 N，并将 (θ) 应用于商模 M/N。）

\begin{enumerate}
\def\labelenumi{(\alph{enumi})}
\setcounter{enumi}{1}
\tightlist
\item
  对于 Z-模 E 为内射模，其充要条件是：对所有 \(x \in E\) 以及每个整数 \(n \neq 0\) ，存在 \(y \in E\) 使得 ny = x。特别地，Z-模 Q 和 Q/Z 是内射的。
\end{enumerate}

\begin{enumerate}
\def\labelenumi{\arabic{enumi}.}
\setcounter{enumi}{11}
\tightlist
\item
  \begin{enumerate}
  \def\labelenumii{(\alph{enumii})}
  \tightlist
  \item
    A-模的一个乘积 \(\prod_{\lambda \in L} E_{\lambda}\) （习题 11）为内射的充要条件是每个 \(E_{\lambda}\) 都是内射的。
  \end{enumerate}
\end{enumerate}

\begin{enumerate}
\def\labelenumi{(\alph{enumi})}
\setcounter{enumi}{1}
\tightlist
\item
  设 K 是一个交换域，L 是一个无限集，A 是乘积环 \(K^{L}\) 。证明 A 是一个内射 A-模（利用 (a)），但 A 的理想 \(\mathfrak{a}\) —— \(K^{L}\) 的因子的直和（典范地等同于 A 的理想）——不是内射 A-模。
\end{enumerate}

¶ 13. (a) 设 A 为一个环，I 为一个内射 A-模（习题 11），使得对于每一个非零单生成 A-模 E，都存在一个从 E 到 I 的非零 A-同态。设 M 为一个 A-模，N 为 M 的一个子模，且 \(\tilde{\mathbf{N}}\) 为由 \(u \in \operatorname{Hom}_{\mathbf{A}}(\mathbf{M}, \mathbf{I})\) （满足 \(u(x) = 0\) 于 N 中）组成的集合。证明对于所有 \(y \notin N\) ，存在一个 \(u \in \tilde{\mathbf{N}}\) ，使得 \(u(y) \neq 0\) 。

\begin{enumerate}
\def\labelenumi{(\alph{enumi})}
\setcounter{enumi}{1}
\item
  设 A 为一个环。对于每个 (A, A)-双模 T 以及每个左（相应地，右）A-模 M， \(\operatorname{Hom}_{\mathbf{A}}(\mathbf{M}, \mathbf{T})\) 典范地具有一个右（相应地，左）A-模结构，它源自 T 上的模结构，并且存在一个典范 A-同态 \(c_{M,T}: M \to \operatorname{Hom}_{\mathbf{A}}(\operatorname{Hom}_{\mathbf{A}}(\mathbf{M}, \mathbf{T}), \mathbf{T})\) ，使得 \(c_{M,T}(x)\) 是 A-同态 \(u \mapsto u(x)\) 。证明：若 I 是一个 (A, A)-双模，且作为左 A-模满足 (a) 中的条件，则典范同态 \(c_{M,I}\) 对每个左 A-模 M 都是单射。
\item
  设 I 是一个 (A, A)-双模，它作为左 A-模和右 A-模都是内射的，并且对于每一个（左或右）单生成 A-模 \(E \neq 0\) ，都存在从 E 到 I 的非零 A-同态。证明在这些条件下，若 P 是左（相应地，右）投射 A-模，则 \(\operatorname{Hom}_{\mathbf{A}}(\mathbf{P}, \mathbf{I})\) 是右（相应地，左）内射 A-模。（假设 P 是左投射 A-模。注意对右 A-模 N 的每个子模 \(N'\) ， \(\operatorname{Hom}_{\mathbf{A}}(\mathbf{N}', \mathbf{I})\) 同构于 \(\operatorname{Hom}_{\mathbf{A}}(\mathbf{N}, \mathbf{I})\) 的一个商模，并将 (b) 应用于右模 N 和左模 P。）
\item
  设 C 为交换环，A 为 C-代数，I 为内射 C-模，使得对每个单生成 C-模 \(E \neq \{0\}\) 存在一个从E到I的非零C-同态。对于每个A-模M， \(\operatorname{Hom}_{\mathbf{C}}(\mathbf{M}, \mathbf{I})\) 是一个A-模；证明对于每个投射A-模P， \(\operatorname{Hom}_{\mathbf{C}}(\mathbf{P}, \mathbf{I})\) 是一个内射A-模（论证与(c)相同）。
\end{enumerate}

\begin{enumerate}
\def\labelenumi{\arabic{enumi}.}
\setcounter{enumi}{13}
\tightlist
\item
  设 A 为一个环。证明：对于每一个 A-模 M，都存在一个内射 A-模 I，使得 M 同构于 I 的一个子模。（取 C = Z，应用习题 13(d)，并利用习题 11(b) 以及每个 A-模都是某个自由 A-模的商模这一事实。）
\end{enumerate}

¶ 15. 一个 A-模同态 \(u: \mathbf{E} \to \mathbf{F}\) 称为本性的，如果它是单射的，并且对于每个子模 \(\mathbf{P} \neq \{0\}\) （ \(\mathbf{F}\) 的）， \(u^{-1}(\mathbf{P}) \neq \{0\}\) ；只需该条件对每个单生成子模 \(\mathbf{P} \neq \{0\}\) （ \(\mathbf{F}\) 的）成立即可。若 \(\mathbf{E}\) 是 \(\mathbf{F}\) 的子模，则称 \(\mathbf{F}\) 为 \(\mathbf{E}\) 的本性扩张，如果典范单射 \(\mathbf{E} \to \mathbf{F}\) 是本性的。

\begin{enumerate}
\def\labelenumi{(\alph{enumi})}
\item
  设 \(u: \mathbf{E} \to \mathbf{F}\) , \(v: \mathbf{F} \to \mathbf{G}\) 是两个 A-模的 A-同态。如果 \(u\) 和 \(v\) 都是本性的，则 \(v \circ u\) 也是本性的。反之，若 \(v \circ u\) 是本性的且 \(v\) 是单射，则 \(u\) 和 \(v\) 都是本性的。
\item
  设 E 是 A-模 F 的一个子模，并设 \((\mathbf{F}_{\lambda})_{\lambda \in L}\) 是 F 的一族子模，包含 E 且其并集为 F。证明：若每个 \(F_{\lambda}\) 都是 E 的本性扩张，则 F 也是。
\item
  设 \((u_{\lambda})_{\lambda \in L}\) 是一族本性同态 \(u_{\lambda} \colon \mathbf{E}_{\lambda} \to \mathbf{F}_{\lambda}\) 。证明同态 \(\bigoplus_{\lambda \in L} u_{\lambda} \colon \bigoplus_{\lambda \in L} \mathbf{E}_{\lambda} \to \bigoplus_{\lambda \in L} \mathbf{F}_{\lambda}\) 是本性的。（先就 \(L\) 只有两个元素的情形证明，然后使用 (b)。）
\item
  \(\mathbf{Z}\) -模 \(\mathbf{Q}\) 是 \(\mathbf{Z}\) 的一个本性扩张，但乘积 \(\mathbf{Q}^{\mathrm{N}}\) 不是 \(\mathbf{Z}^{\mathrm{N}}\) 的本性扩张。
\item
  设 E 是模 F 的一个子模。证明存在一个子模
\end{enumerate}

F 的一个子模 Q，使得 \(Q \cap E = \{0\}\) ，并且典范同态 \(F \to F/Q\) 在 E 上的限制是本性的。

\begin{enumerate}
\def\labelenumi{\arabic{enumi}.}
\setcounter{enumi}{15}
\tightlist
\item
  A-模 F 的子模 E 称为关于 F（或在 F 中）不可约的，如果 \(E \neq F\) 且 E 不是 F 的两个不同于 E 的子模的交集。
\end{enumerate}

\begin{enumerate}
\def\labelenumi{(\alph{enumi})}
\item
  A-模 F 是其每个子模 \(\neq \{0\}\) 的本性扩张的充要条件是： \(\{0\}\) 关于 F 是不可约的。
\item
  设 \((\mathbf{E}_{\lambda})_{\lambda \in L}\) 是 A-模 \(\mathbf{F}\) 的一族子模。称 \((\mathbf{E}_{\lambda})_{\lambda \in L}\) 为 \(\{0\}\) 在 \(\mathbf{F}\) 中的一个既约不可约分解，如果每个 \(\mathbf{F}_{\lambda}\) 都在 \(\mathbf{F}\) 中不可约，诸 \(\mathbf{E}_{\lambda}\) 的交集化为 0，并且没有哪个 \(\mathbf{E}_{\lambda}\) 包含诸 \(\mathbf{E}_{\mu}\) （ \(\mu \neq \lambda\) ）的交集。在此情形下，证明 \(\mathbf{F}\) 到 \(\bigoplus_{\lambda \in L} (\mathbf{F} / \mathbf{E}_{\lambda})\) 中的典范同态是本性的（若 \(\mathbf{E}_{\lambda}' = \bigcap_{\mu \neq \lambda} \mathbf{E}_{\mu}\) ，注意像 \(\mathbf{F}_{\lambda}'\) （ \(\mathbf{E}_{\lambda}'\) 在 \(\mathbf{F} / \mathbf{E}_{\lambda}\) 中的）非零，且 \(\mathbf{F}\) 在 \(\bigoplus_{\lambda \in L} (\mathbf{F} / \mathbf{E}_{\lambda})\) 中的像包含 \(\bigoplus_{\lambda \in L} \mathbf{F}_{\lambda}'\) ；然后应用习题 15(c)）。反之，若 \((\mathbf{E}_{\lambda})_{\lambda \in L}\) 是 \(\mathbf{F}\) 的子模的一个族，在 \(\mathbf{F}\) 中不可约，且典范同态 \(\mathbf{F} \to \bigoplus_{\lambda \in L} (\mathbf{F} / \mathbf{E}_{\lambda})\) 是本性的，则族 \((\mathbf{E}_{\lambda})\) 是 \(\{0\}\) 在 \(\mathbf{F}\) 中的一个既约不可约分解。
\end{enumerate}

¶ 17. (a) 设 M, N, M', N' 是模 E 的四个子模，使得 M ∩ N = M' ∩ N' = P。证明

\[
\mathbf {M} = (\mathbf {M} + (\mathbf {N} \cap \mathbf {M} ^ {\prime})) \cap (\mathbf {M} + (\mathbf {N} \cap \mathbf {N} ^ {\prime})).
\]

由此推出，如果 \(\mathbf{M}\) 在 \(\mathbf{E}\) 中不可约，则必然有 \(\mathbf{P} = \mathbf{M}' \cap \mathbf{N}\) 或 \(\mathbf{P} = \mathbf{N}' \cap \mathbf{N}\) 。

\begin{enumerate}
\def\labelenumi{(\alph{enumi})}
\setcounter{enumi}{1}
\tightlist
\item
  设 \((\mathrm{N}_{i})_{1\leqslant i\leqslant n}\) 是 E 中不可约子模的有限族，M 为其交集。证明若 \((\mathrm{P}_{j})_{1\leqslant j\leqslant m}\) 是 E 的子模的有限族，使得 \(M=\bigcap_{j}P_{j}\) ，那么对于每个满足 \(1\leqslant i\leqslant n\) 的指标 i，存在一个指标 \(\phi(i)\) ，使得 \(1\leqslant\phi(i)\leqslant m\) ，并且使得，如果我们写 \(N_{i}^{\prime}=\bigcap_{k\neq i}N_{k}\) ，则 \(M=P_{\phi(i)}\cap N_{i}^{\prime}\) （利用 (a)）。由此推出，如果没有哪个 \(P_{j}\) 包含诸 \(P_{k}\) （下标 \(\neq j\) 者）的交集，则 \(m\leqslant n\) （依次将 \(N_{i}\) 替换为合适的 \(P_{j}\) ，利用上述结果）。由此得出结论：M 在 E 中的两个既约不可约分解，若其中一个是有限的，则它们必然具有相同的项数（参见习题 23）。
\end{enumerate}

¶ 18. (a) 为使一个 A-模 I 是内射的，证明其必要且充分条件是 I 不具有任何不同于自身的本性扩张（证明此条件蕴含习题 11(a) 的条件 (δ)，论证方式与该习题中 (θ) 蕴含 (δ) 的证明相同）。

\begin{enumerate}
\def\labelenumi{(\alph{enumi})}
\setcounter{enumi}{1}
\item
  设 I 为内射 A-模。对于 I 的子模 E 而言，E 为内射模当且仅当 E 不具有任何不同于自身且包含在 I 中的本性扩张（利用习题 15(e)），证明此条件蕴含 E 是 I 的直因子）。
\item
  设 I 是内射 A-模，E 是 I 的子模，且 \(h: \mathbf{E} \to \mathbf{F}\) 是一个本性同态。证明存在一个单同态 \(j: \mathbf{F} \to \mathbf{I}\) ，使得 \(j \circ h\) 是 E 到 I 的典范单射。
\item
  设 I 为 A-模，E 为 I 的子模。证明以下性质等价：
\end{enumerate}

(α) I 是内射的，并且是 E 的一个本性扩张。

(β) I 是内射的，并且对于每个单同态 \(j\) （从 \(E\) 到一个内射 A-模 \(J\) 中的），存在一个从 \(I\) 到 \(J\) 的单同态，它扩张 \(j\) 。

(γ) I 是内射的，并且是 I 的包含 E 的最小内射子模。

(δ) I 是 E 的本性扩张，并且对于每个本性同态 \(h: E \to F\) ，存在一个单同态 \(j: F \to I\) ，使得 \(j \circ h\) 是 E 到 I 的典范单射。

（利用 (c) 和 (a)。）当这些等价条件成立时，I 称为 E 的内射包络；此时 I 也是 I 的包含 E 的每个子模的内射包络。

\begin{enumerate}
\def\labelenumi{(\alph{enumi})}
\setcounter{enumi}{4}
\item
  设 I 是内射 A-模，E 是 I 的子模。证明：包含在 I 中的 E 的本性扩张所成的集合的每个极大元都是 E 的一个内射包络（利用 (b)）。特别地，每个 A-模都具有内射包络（参见习题 14）。
\item
  若 I, I' 是同一个 A-模 E 的两个内射包络，证明存在从 I 到 I' 的同构，且该同构保持 E 的元素不变。
\end{enumerate}

¶ 19. (a) 设 E, F 是两个 A-模，M 是 E ⊕ F 的一个内射子模。设 I 是 M ∩ E 在 M 中的内射包络，J 是 I 在 M 中的互补子模。证明 E ⊕ F 到 E（相应地，F）上的典范投影在 I（相应地，J）上的限制是单射（将 M ∩ E → I 的单射的一个扩张 E → I 与 I 到 E 上的投影复合）。

\begin{enumerate}
\def\labelenumi{(\alph{enumi})}
\setcounter{enumi}{1}
\tightlist
\item
  设 E 是 A-模，M 是 E 的子模，且在 E 的内射子模集合中是极大的；M 在 E 中具有一个互补子模 N，它不包含任何内射子模 \(\neq\{0\}\) 。证明：对于 E 的每个内射子模 P，P 在 E 到 M 上的投影（在 E 的直和分解 \(M \oplus N\) 下）下的像是 \(P \cap M\) 在 M 中的一个内射包络（利用 (a)）。若 \(M'\) 是 E 的另一个子模，且在 E 的内射子模集合中是极大的，证明存在 E 的一个自同构将 M 变为 \(M'\) 并保持 N 的元素不变。
\end{enumerate}

¶ 20. (a) 设 A 是一个环，具有一个左分式域 K（I，§ 9，习题 15）。证明 K 看作左 A-模时是 A \(_{s}\) 的内射包络。（应用习题 18(d) 的判据 (α) 和习题 11(a) 的判据 (δ)，注意到若 f: a → K 是 A 的左理想 a 到 K 中的 A-线性映射，则元素 \(x^{-1}f(x)\) （对于 \(x \in \mathfrak{a}, x \neq 0\) ，逆元在 K 中取）不依赖于 \(x \in \mathfrak{a}\) 。）

\begin{enumerate}
\def\labelenumi{(\alph{enumi})}
\setcounter{enumi}{1}
\tightlist
\item
  设 \(U_{p}\) 为 Q/Z 的子 Z-模，由形如 \(k/p^{n}\) 的有理数的典范像组成，其中 p 是给定的素数， \(k \in Z\) 和 \(n \in N\) （II，§1，习题 12(d)）。证明 \(U_{p}\) 是每个 Z-模 \(p^{-n}Z/Z\) （同构于 \(Z/p^{n}Z\) 的）的内射包络。证明存在 \(U_{p}\) 的自同构，它们不同于恒等自同构，但保持某个子模 \(p^{-n}Z/Z\) 的元素不变。
\end{enumerate}

¶ 21. (a) 设 \((\mathbf{E}_j)_{j \in J}\) 是 A-模的一个有限族，且对所有 \(j \in J\) ，设 \(I_j\) 是 \(E_j\) 的一个内射包络；证明 \(\bigoplus_{j \in J} I_j\) 是 \(\bigoplus_{j \in J} E_j\) 的一个内射包络。

\begin{enumerate}
\def\labelenumi{(\alph{enumi})}
\setcounter{enumi}{1}
\tightlist
\item
  一个 A-模 M 称为不可分解的，如果它不同于 \(\{0\}\) 并且它不是两个不同于 \(\{0\}\) 的子模的直和。证明：若 I 是一个内射 A-模，则以下性质等价：
\end{enumerate}

\((\alpha)\) \(\{0\}\) 是 I 中的一个不可约子模。

(β) I 是不可分解的。

(γ) I 是其每个子模 \(\neq0\) 的内射包络。（利用 (a) 和习题 18 (e)。）

(δ) I 同构于内射包络 \(\mathrm{I}(\mathrm{A}/\mathrm{q})\) ，其中 q 是 A 的一个不可约左理想。

此外，当这种情况成立时，对 I 中所有 \(x \neq 0\) ， \(\operatorname{Ann}(x)\) 是 A 的一个不可约左理想，且 I 同构于 \(\operatorname{I}(\mathbf{A} / \operatorname{Ann}(x))\) 。

由此推出，A-模 E 的内射包络是不可分解的，其充分必要条件为 \(\{0\}\) 是 E 的一个不可约子模。

\begin{enumerate}
\def\labelenumi{(\alph{enumi})}
\setcounter{enumi}{2}
\item
  给出一个不可分解 Z-模 E 的例子，使得其内射包络 I(E) 是可分解的（参见 VII，§3，习题 5）。
\item
  设 E 是一个 A-模，I 是 E 的一个内射包络，且 \(I = \bigoplus_{\lambda \in L} I_{\lambda}\) 是 I 作为有限个不可分解内射子模的直和的一个分解；对所有 \(\lambda \in L\) ，设 \(J_{\lambda} = \bigoplus_{\mu \neq \lambda} I_{\mu}\) ，并设 \(N_{\lambda} = J_{\lambda} \cap E\) 。证明 \((\mathrm{N}_{\lambda})_{\lambda \in L}\) 是 \(\{0\}\) 在 E 中的一个既约不可约分解（习题 16 (b)），并且 \(I_{\lambda}\) 是 \(E/N_{\lambda}\) 的一个内射包络。反之， \(\{0\}\) 在 E 中的每一个既约不可约分解都可以由上述过程唯一地（在同构意义下）得到（使用习题 16 (b)）。
\end{enumerate}

\begin{enumerate}
\def\labelenumi{\arabic{enumi}.}
\setcounter{enumi}{21}
\tightlist
\item
  设 I 是一个内射 A-模。若 I 是不可分解的，则 I 的每一个单射自同态都是 I 的自同构。由此推出，I 不可分解的充分必要条件是环 End(I) 中所有不可逆元素构成该环的一个双边理想。
\end{enumerate}

¶ 23. 设 M 是一个 A-模，它是一个子模族（有限或非有限） \((\mathbf{M}_{\lambda})_{\lambda \in L}\) 的直和，使得对所有 \(\lambda \in L\) ， \(\operatorname{End}(\mathbf{M}_{\lambda})\) 的不可逆元素构成该环的一个双边理想（这意味着 \(\mathbf{M}_{\lambda}\) 是不可分解的）。

\begin{enumerate}
\def\labelenumi{(\alph{enumi})}
\item
  设 \(f, g\) 是 \(\mathbf{M}\) 的两个自同态，满足 \(f + g = 1_{\mathbf{M}}\) 。证明：对于互异指标的每个有限序列 \((\lambda_k)_{1 \leqslant k \leqslant s}\) （在 \(\mathbf{L}\) 中取的），存在一个子模族 \(\mathbf{N}_k\) ( \(1 \leqslant k \leqslant s\) )（ \(\mathbf{M}\) 的），使得对每个 \(k\) ，自同态 \(f, g\) 之一限制在 \(\mathbf{M}_{\lambda_{k}}\) 上时是该子模到 \(\mathbf{N}_k\) 上的一个同构，并且 \(\mathbf{M}\) 是诸 \(\mathbf{N}_k\) ( \(1 \leqslant k \leqslant s\) ) 与诸 \(\mathbf{M}_{\lambda}\) （ \(\lambda\) 不同于 \(\lambda_k\) 者）的直和。（化简到 \(s = 1\) 的情形；若 \(\pi_{\lambda}\) 是 \(\mathbf{M}\) 到 \(\mathbf{M}_{\lambda}\) 上的典范投影——对应于 \(\mathbf{M}_{\lambda}\) 与 \(\bigoplus_{\mu \neq \lambda} \mathbf{M}_{\mu}\) 的直和分解——注意 \(\pi_{\lambda} \circ f\) 或 \(\pi_{\lambda} \circ g\) 之一限制在 \(\mathbf{M}_{\lambda}\) 上时是 \(\mathbf{M}_{\lambda}\) 的一个自同构。）
\item
  设 \(f\) 是 \(\mathbf{M}\) 的一个幂等自同态。证明存在至少一个指标 \(\lambda \in \mathbf{L}\) ，使得 \(f\) 在 \(\mathbf{M}_{\lambda}\) 上的限制是 \(\mathbf{M}_{\lambda}\) 到 \(f(\mathbf{M}_{\lambda})\) 上的一个同构；此外 \(f(\mathbf{M}_{\lambda})\) 是 \(\mathbf{M}\) 的一个直因子。（若 \(g = 1_{\mathbf{M}} - f\) ，注意存在一个有限的指标序列 \((\lambda_k)_{1 \leqslant k \leqslant s}\) ，使得 \(\operatorname{Ker}(g)\) 与诸 \(\mathbf{M}_{\lambda_k}\) 的直和的交 \(\neq \{0\}\) ，并利用 (a)。）
\item
  从 (b) 推出，M 的每个不可分解直因子都同构于某个 \(M_{\lambda}\) 。
\item
  设 \((\mathrm{N}_{\kappa})_{\kappa\in K}\) 是 M 的一族不可分解子模，M 是它们的直和；由 (c)，每个 \(N_{\kappa}\) 同构于某个 \(M_{\lambda}\) ，反之亦然。设 T 为 M 的不可分解子模的类（在同构关系下）的集合，使得某个 \(M_{\lambda}\) （或某个 \(N_{\kappa}\) ）属于这些类之一。对所有 \(t\in T\) ，设 \(\mathrm{R}(t)\) （相应地 \(\mathrm{S}(t)\) ）为 \(\lambda\in L\) （相应地 \(\kappa\in K\) ）中使 \(M_{\lambda}\in t\) （相应地 \(N_{\kappa}\in t\) ）者组成的集合。证明：对所有 \(t\in T\) , \(\mathrm{Card}(\mathrm{R}(t))=\mathrm{Card}(\mathrm{S}(t))\) 。（在 (a) 的记号中，令 \(\mathrm{J}(\kappa)\) 为由满足如下条件的 \(\lambda\in L\) 组成的集合： \(\pi_{\lambda}\) 在 \(N_{\kappa}\) 上的限制是 \(N_{\kappa}\) 到 \(M_{\lambda}\) 上的一个同构；证明 \(\mathrm{J}(\kappa)\) 是有限的，并且诸 \(\mathrm{J}(\kappa)\) 覆盖 \(\mathrm{R}(t)\) （当 \(\kappa\) 遍历 \(\mathrm{S}(t)\) 时）。由此推出， \(\mathrm{Card}(\mathrm{S}(t))\leqslant\mathrm{Card}(\mathrm{R}(t))\) （当 \(\mathrm{R}(t)\) 无限时）。当 \(\mathrm{R}(t)\) 有限时，借助 (a) 证明 M 是诸 \(M_{\lambda}\) （ \(\lambda\notin R(t)\) ）与 \((\mathrm{N}_{\kappa})_{\kappa\in\mathrm{S}(t)}\) 的一个基数为 \(\mathrm{R}(t)\) 的子族的直和。
\item
  由 (d) 以及习题 22 和习题 21(d) 推出，若 \((\mathbf{E}_{\lambda})_{\lambda \in L}\) 和 \((\mathbf{E}_{\kappa}^{\prime})_{\kappa \in K}\) 是 \(\{0\}\) 在模 F 中的两个既约不可约分解，则 L 与 K 是等势的。
\end{enumerate}

\begin{enumerate}
\def\labelenumi{\arabic{enumi}.}
\setcounter{enumi}{23}
\item
  设 M 是一个 A-模，I 是它的内射包络（习题 18）， \(\mathfrak{a}\) 是 A 的一个双边理想，Q 为 I 的由 \(z \in I\) （满足 \(\mathfrak{a}z = \{0\}\) ）组成的子模；Q 具有一个自然的 \((\mathrm{A}/\mathfrak{a})\) -模结构。设 N 为 M 的由 \(x \in M\) （满足 \(\mathfrak{a}x = \{0\}\) ）组成的子模；若将 N 视为一个 \((\mathrm{A}/\mathfrak{a})\) -模，证明 Q 同构于 N 的一个内射包络。
\item
  \begin{enumerate}
  \def\labelenumii{(\alph{enumii})}
  \tightlist
  \item
    要使 A-模 \(\mathbf{Q}\) 是内射的，只需：对每个投射 A-模 P 及 P 的每个子模 \(P'\) ，每个从 \(P'\) 到 Q 的线性映射都能扩张为从 P 到 Q 的线性映射（利用每个 A-模都是某个投射 A-模的商模这一事实）。
  \end{enumerate}
\end{enumerate}

\begin{enumerate}
\def\labelenumi{(\alph{enumi})}
\setcounter{enumi}{1}
\tightlist
\item
  要使 A-模 P 是投射的，只需：对每个内射 A-模 Q 及 Q 的每个商模 \(Q''\) ，每个从 P 到 \(Q''\) 中的同态都具有 \(\phi \circ u\) 的形式，其中 \(\phi: Q \to Q''\) 是典范同态，而 u 是从 P 到 Q 中的同态（利用每个 A-模都是某个内射 A-模的子模这一事实）。
\end{enumerate}

\begin{enumerate}
\def\labelenumi{\arabic{enumi}.}
\setcounter{enumi}{25}
\tightlist
\item
  对于每个 Z-模 G 和每个整数 n \textgreater{} 0，令 \(_{n}G\) 表示 G 的自同态 \(x \mapsto nx\) 的核。若 \(0 \to G' \to G \to G'' \to 0\) 是正合序列，定义典范同态 \(d:_{n}G'' \to G'/nG'\) 使得序列
\end{enumerate}

\[
0 \longrightarrow_ {n} G ^ {\prime} \longrightarrow_ {n} G ^ {\prime \prime} \stackrel {d} {\longrightarrow} G ^ {\prime} / n G ^ {\prime} \longrightarrow G / n G \longrightarrow G ^ {\prime \prime} / n G ^ {\prime \prime} \longrightarrow 0
\]

正合。若 \(_n\mathrm{G}\) 和 \(\mathrm{G} / n\mathrm{G}\) 是有限的，我们记

\[
h _ {n} (\mathrm{G}) = \operatorname{Card} (\mathrm{G} / n \mathrm{G}) - \operatorname{Card} (_ {n} \mathrm{G});
\]

则证明 \(_{n}G'\) , \(_{n}G''\) , \(G'/nG'\) , \(G''/nG''\) 是有限的，并且

\[
h _ {n} (\mathrm{G}) = h _ {n} \left(\mathrm{G} ^ {\prime}\right) + h _ {n} \left(\mathrm{G} ^ {\prime \prime}\right).
\]

